% Part: first-order-logic
% Chapter: syntax-and-semantics
% Section: terms-formulas

\documentclass[../../../include/open-logic-section]{subfiles}

\begin{document}

\olfileid{fol}{syn}{frm}

\olsection{પદો અને \printtoken{P}{formula}}

એક વાર પ્રથમ-ક્રમ ભાષા~$\Lang L$ આપવામાં આવે પછી, $\Lang L$ના પાયાના
સંકેતભંડોળમાંથી બનતી અભિવ્યક્તિઓ વ્યાખ્યાયિત કરી શકીએ છીએ. તેમાં ખાસ
કરીને \emph{પદો} અને \emph{!!{formula}s}નો સમાવેશ થાય છે.

\begin{defn}[પદો]
\ollabel{defn:terms}
$\Lang L$નાં \emph{પદો}નો ગણ~$\Trm[L]$ નીચે પ્રમાણે અનુમાનાત્મક રીતે
વ્યાખ્યાયિત થાય છે:
\begin{enumerate}
\item દરેક !!{variable} પદ છે.
\item $\Lang L$નું દરેક !!{constant} પદ છે.
\item જો $f$ કોઈ $n$-સ્થાની !!{function} હોય અને $t_1$, \dots, $t_n$
  પદો હોય, તો $\Atom{f}{t_1, \ldots, t_n}$ પદ છે.
\tagitem{limitClause}{
  બીજું કશું પદ નથી.}{}
\end{enumerate}
!!{variable}sમાંથી એક પણ ન ધરાવતું પદ \emph{બંધ પદ} કહેવાય છે.
\end{defn}

\begin{explain}
ભાષા અને પદોના આપણા નિર્દેશમાં !!{constant}sને સંકેતોની અલગ કોટિ
તરીકે મૂક્યાં છે, પરંતુ તેમને શૂન્ય-સ્થાની !!{function}sમાં સામેલ કરી
શકાય. ત્યારે પદોની વ્યાખ્યાની બીજી કલમની જરૂર રહે નહિ. માત્ર એ
સમજવાનું રહે કે $n=0$ હોય ત્યારે $\Atom{f}{t_1, \ldots, t_n}$નો અર્થ
એકલું~$f$ છે.
\end{explain}

\begin{defn}[સૂત્રો]
\ollabel{defn:formulas}
ભાષા~$\Lang L$નાં \emph{!!{formula}s}નો ગણ~$\Frm[L]$ નીચે પ્રમાણે
અનુમાનાત્મક રીતે વ્યાખ્યાયિત થાય છે:
\begin{enumerate}
\tagitem{prvFalse}{$\lfalse$ આણ્વિક !!{formula} છે.}{}

\tagitem{prvTrue}{$\ltrue$ આણ્વિક !!{formula} છે.}{}

\item જો $R$, $\Lang L$નું $n$-સ્થાની !!{predicate} હોય અને $t_1$,
  \dots, $t_n$, $\Lang L$નાં પદો હોય, તો
  $\Atom{R}{t_1,\ldots,t_n}$ આણ્વિક !!{formula} છે.

\item જો $t_1$ અને $t_2$, $\Lang L$નાં પદો હોય, તો
  $\Atom{\eq}{t_1,t_2}$ આણ્વિક !!{formula} છે.

\tagitem{prvNot}{જો $!A$ !!a{formula} હોય, તો $\lnot !A$ પણ
  !!a{formula} છે.}{}

\tagitem{prvAnd}{જો $!A$ અને $!B$ !!{formula}s હોય, તો $(!A \land
  !B)$ પણ !!a{formula} છે.}{}

\tagitem{prvOr}{જો $!A$ અને $!B$ !!{formula}s હોય, તો $(!A \lor !B)$
  પણ !!a{formula} છે.}{}

\tagitem{prvIf}{જો $!A$ અને $!B$ !!{formula}s હોય, તો $(!A \lif !B)$
  પણ !!a{formula} છે.}{}

\tagitem{prvIff}{જો $!A$ અને $!B$ !!{formula}s હોય, તો $(!A \liff !B)$
  પણ !!a{formula} છે.}{}

\tagitem{prvAll}{જો $!A$ !!a{formula} અને $x$ !!a{variable} હોય,
  તો $\lforall[x][!A]$ પણ !!a{formula} છે.}{}

\tagitem{prvEx}{જો $!A$ !!a{formula} અને $x$ !!a{variable} હોય,
  તો $\lexists[x][!A]$ પણ !!a{formula} છે.}{}

\tagitem{limitClause}{બીજું કશું !!a{formula} નથી.}{}
\end{enumerate}
\end{defn}

\begin{explain}
પદોના ગણની અને !!{formula}sના ગણની વ્યાખ્યાઓ \emph{અનુમાનાત્મક
વ્યાખ્યાઓ} છે. મૂળભૂત રીતે, આપણે !!{formula}sનો ગણ અનંત સંખ્યામાં
તબક્કાઓમાં રચીએ છીએ. પહેલા તબક્કામાં બધાં આણ્વિક સૂત્રોને સૂત્ર જાહેર
કરીએ છીએ; આ વ્યાખ્યાના પહેલા થોડા કિસ્સાઓ, એટલે કે
\iftag{prvTrue}{$\ltrue$, }{}%
\iftag{prvFalse}{$\lfalse$, }{}%
$\Atom{R}{t_1,\dots,t_n}$ અને $\Atom{\eq}{t_1,t_2}$ના કિસ્સાઓને
અનુરૂપ છે. તેથી ``આણ્વિક !!{formula}''નો અર્થ આમાંના કોઈ પણ સ્વરૂપનું
!!{formula} થાય છે.

વ્યાખ્યાના બીજા કિસ્સાઓ અગાઉથી રચેલાં !!{formula}sમાંથી નવાં
!!{formula}s રચવાના નિયમો આપે છે. બીજા તબક્કામાં આ નિયમો વડે આણ્વિક
!!{formula}sમાંથી !!{formula}s રચી શકીએ છીએ. ત્રીજા તબક્કામાં આણ્વિક
સૂત્રો અને બીજા તબક્કામાં મળેલાં સૂત્રોમાંથી નવાં સૂત્રો રચીએ છીએ,
અને આમ આગળ ચાલે છે. આવા કોઈ તબક્કે છેવટે રચાતી દરેક વસ્તુ, અને માત્ર
એવી વસ્તુ, !!{formula} છે.
\end{explain}

પરંપરા મુજબ આપણે $\eq$ને તેની દલીલોની વચ્ચે લખીને કૌંસ છોડી દઈએ
છીએ: $\eq[t_1][t_2]$ એ $\Atom{\eq}{t_1,t_2}$નું સંક્ષિપ્ત રૂપ છે.
વળી, $\lnot \Atom{\eq}{t_1,t_2}$ને $\eq/[t_1][t_2]$ તરીકે ટૂંકાવીએ
છીએ. દ્વિસ્થાની સંયોજક~$\ast$ વડે $!B$, $!C$માંથી રચેલું સૂત્ર
$(!B \ast !C)$ લખતી વખતે આપણે ઘણી વાર સૌથી બહારની કૌંસજોડી છોડી
દઈને માત્ર~$!B \ast !C$ લખીશું.

\begin{intro}
કેટલાંક તર્કશાસ્ત્રનાં પુસ્તકોમાં
\iftag{prvEx}{$\lexists[x][!A]$ }{}%
\iftag{notprvEx,notprvAll}{}{અને }%
\iftag{prvAll}{$\lforall[x][!A]$ }{}%
ને
\iftag{notprvEx,notprvAll}{!!a{formula}}{!!{formula}s}
ગણવા માટે !!{variable}~$x$નું~$!A$માં આવવું જરૂરી રાખે છે. આ શરત
ન રાખવાથી કશું ખોટું થતું નથી અને કામ સરળ બને છે.
\end{intro}

\begin{tagblock}{defNot,defOr,defAnd,defIf,defIff,defTrue,defFalse,defEx,defAll}
\begin{defn}
વ્યાખ્યાયિત કારકો વાપરીને રચેલાં સૂત્રો નીચે પ્રમાણે સમજવાનાં છે:

\begin{tagenumerate}{defTrue,defFalse,defNot,defOr,defAnd,defIf,defIff,defEx,defAll}

\tagitem{defTrue}{$\ltrue$ એ
  \iftag{prvFalse}{$\lnot\lfalse$}{$(!A \lor \lnot !A)$, જ્યાં
    $!A$ કોઈ નિશ્ચિત આણ્વિક !!{formula} છે}, તેનું સંક્ષિપ્ત રૂપ છે.}{}

\tagitem{defFalse}{$\lfalse$ એ
  \iftag{prvTrue}{$\lnot\ltrue$}{$(!A \land \lnot !A)$, જ્યાં
    $!A$ કોઈ નિશ્ચિત આણ્વિક !!{formula} છે}, તેનું સંક્ષિપ્ત રૂપ છે.}{}

\tagitem{defNot}{$\lnot !A$ એ $!A \lif \lfalse$નું સંક્ષિપ્ત રૂપ છે.}{}

\tagitem{defOr}{$!A \lor !B$ એ
  \iftag{prvAnd}{$\lnot(\lnot !A \land \lnot !B)$}{$\lnot !A \lif
    !B$}નું સંક્ષિપ્ત રૂપ છે.}{}

\tagitem{defAnd}{$!A \land !B$ એ
  \iftag{prvOr}{$\lnot(\lnot !A \lor \lnot !B)$}{$\lnot (!A \lif
    \lnot !B)$}નું સંક્ષિપ્ત રૂપ છે.}{}

\tagitem{defIf}{$!A \lif !B$ એ
  \iftag{prvOr}{$\lnot !A \lor !B$}{$\lnot (!A \land \lnot !B)$}નું
  સંક્ષિપ્ત રૂપ છે.}{}

\tagitem{defIff}{$!A \liff !B$ એ $(!A \lif !B) \land (!B
  \lif !A)$નું સંક્ષિપ્ત રૂપ છે.}{}

\tagitem{defAll}{$\lforall[x][!A]$ એ $\lnot\lexists[x][\lnot !A]$નું
  સંક્ષિપ્ત રૂપ છે.}{}

\tagitem{defEx}{$\lexists[x][!A]$ એ $\lnot\lforall[x][\lnot !A]$નું
  સંક્ષિપ્ત રૂપ છે.}{}
\end{tagenumerate}
\end{defn}
\end{tagblock}
\sourcecorrection{OLSYN-002}{મૂળની \texttt{terms-formulas.tex}માં
વ્યાખ્યાયિત પ્રેરણના વિયોજનવાળા વિકલ્પને
$\lnot !A\lor !B)$ લખતાં એક વધારાનું અંતકૌંસ આવે છે. અહીં
અર્થાત્ બનતું સુઘડ સૂત્ર $\lnot !A\lor !B$ પુનઃસ્થાપિત કર્યું છે.}

ચોક્કસ ઉપયોગ માટેની ભાષામાં કામ કરીએ ત્યારે દ્વિસ્થાની
!!{predicate}s અને !!{function}sને સંબંધિત પદોની વચ્ચે લખીએ છીએ;
દાખલા તરીકે, અંકગણિતની ભાષામાં $t_1<t_2$ અને $(t_1+t_2)$ તથા
ગણસિદ્ધાંતની ભાષામાં $t_1\in t_2$. અંકગણિતની ભાષામાં ઉત્તરગામી
વિધેય પરંપરા મુજબ તેની દલીલની \emph{પછી} પણ લખાય છે:~$t'$. જોકે
સત્તાવાર રીતે આ અનુક્રમે $\Atom{\Obj A^2_0}{t_1,t_2}$,
$\Obj f^2_0(t_1,t_2)$, $\Atom{\Obj A^2_0}{t_1,t_2}$ અને
$f^1_0(t)$નાં માત્ર પરંપરાગત સંક્ષિપ્ત રૂપો છે.

\begin{defn}[વાક્યરચનાત્મક અભિન્નતા]
$\ident$ સંકેત, સંકેતોની શ્રેણીઓ વચ્ચેની વાક્યરચનાત્મક અભિન્નતા
દર્શાવે છે; એટલે કે, $!A \ident !B$ ત્યારે જ જ્યારે $!A$ અને $!B$
સમાન લંબાઈ ધરાવતી સંકેતશ્રેણીઓ હોય અને દરેક સ્થાને બંનેમાં એક જ
સંકેત હોય.
\end{defn}

$\ident$ સંકેતની બંને બાજુએ જોડાણથી મળેલી શ્રેણીઓ પણ આવી શકે છે.
દાખલા તરીકે, $!A \ident (!B \lor !C)$નો અર્થ આ છે: $!A$ની
સંકેતશ્રેણી એ આરંભકૌંસ, શ્રેણી~$!B$, સંકેત~$\lor$, શ્રેણી~$!C$ અને
અંતકૌંસને આ જ ક્રમમાં જોડવાથી મળતી શ્રેણી છે. આમ હોય તો આપણને ખબર
પડે છે કે $!A$નો પહેલો સંકેત આરંભકૌંસ છે, $!A$માં $!B$ ઉપશ્રેણી
તરીકે આવે છે (બીજા સંકેતથી શરૂ થઈને), તે ઉપશ્રેણી પછી~$\lor$ આવે છે,
વગેરે.

પદો અને !!{formula}s પાયાના ઘટકોમાંથી અનુમાનાત્મક વ્યાખ્યાઓ વડે
રચાય છે, તેથી તેમના વિશે વાતો સાબિત કરવા આપણે નીચેના અનુમાનના
સિદ્ધાંતો વાપરી શકીએ છીએ.

\begin{lem}[\emph{પદો પર અનુમાનનો સિદ્ધાંત}]
    \ollabel{lem:trmind}
    $\Lang L$ પ્રથમ-ક્રમ ભાષા હોય. કોઈ ગુણધર્મ~$P$ એવો હોય કે
    %
    \begin{enumerate}
        \item તે દરેક !!{variable}~$v$ માટે હોય,
        %
        \item તે $\Lang L$ના દરેક !!{constant}~$a$ માટે હોય, અને
        %
        \item જ્યારે તે $t_1$,~\dots, $t_n$ માટે હોય અને $f$,
        $\Lang L$નું $n$-સ્થાની !!{function} હોય ત્યારે તે
        $f(t_1,\dotsc,t_n)$ માટે પણ હોય
    \end{enumerate}
    ($t_1$,~\dots, $t_n$, $\Lang{L}$નાં પદો છે એમ ધારીને),
    તો $P$, $\Trm[L]$ના દરેક પદ માટે હોય છે.
\end{lem}

\begin{prob}
    \olref[fol][syn][frm]{lem:trmind} સાબિત કરો.
\end{prob}

\begin{lem}[\emph{!!{formula}s પર અનુમાનનો સિદ્ધાંત}]
    \ollabel{thm:frmind}
    $\Lang L$ પ્રથમ-ક્રમ ભાષા હોય. કોઈ ગુણધર્મ~$P$ બધાં આણ્વિક
    !!{formula}s માટે હોય અને એવો હોય કે
    %
    \begin{enumerate}
        \tagitem{prvNot}{જ્યારે તે~$!A$ માટે હોય ત્યારે તે
            $\lnot !A$ માટે પણ હોય;}{}
        \tagitem{prvAnd}{જ્યારે તે $!A$ અને~$!B$ માટે હોય ત્યારે તે
            $(!A \land !B)$ માટે પણ હોય;}{}
        \tagitem{prvOr}{જ્યારે તે $!A$ અને~$!B$ માટે હોય ત્યારે તે
            $(!A \lor !B)$ માટે પણ હોય;}{}
        \tagitem{prvIf}{જ્યારે તે $!A$ અને~$!B$ માટે હોય ત્યારે તે
            $(!A \lif !B)$ માટે પણ હોય;}{}
        \tagitem{prvIff}{જ્યારે તે $!A$ અને~$!B$ માટે હોય ત્યારે તે
            $(!A \liff !B)$ માટે પણ હોય;}{}
        \tagitem{prvEx}{જ્યારે તે~$!A$ માટે હોય ત્યારે તે
            $\lexists[x][!A]$ માટે પણ હોય;}{}
        \tagitem{prvAll}{જ્યારે તે~$!A$ માટે હોય ત્યારે તે
            $\lforall[x][!A]$ માટે પણ હોય;}{}
  \end{enumerate}
  ($!A$ અને $!B$, $\Lang{L}$નાં !!{formula}s છે એમ ધારીને),
  તો $P$, $\Frm[L]$નાં બધાં સૂત્રો માટે હોય છે.
\end{lem}

\end{document}
